\documentclass[11pt,a4paper]{article}
% Preâmbulo comum. Não duplicar pacotes nos módulos.
\usepackage[utf8]{inputenc}
\usepackage[T1]{fontenc}
\usepackage[brazil]{babel}
\usepackage{amsmath,amssymb,amsthm,mathtools}
\usepackage{geometry}
\geometry{a4paper,margin=2.3cm}
\usepackage{microtype}
\usepackage{enumitem}
\usepackage{graphicx}
\usepackage{xcolor}
\usepackage{booktabs}
\usepackage{hyperref}
\hypersetup{colorlinks=true,linkcolor=blue!40!black,urlcolor=blue!50!black}

\setcounter{tocdepth}{2}
\numberwithin{equation}{section}

\theoremstyle{plain}
\newtheorem{teorema}{Teorema}[section]
\newtheorem{proposicao}[teorema]{Proposição}
\newtheorem{lema}[teorema]{Lema}
\newtheorem{corolario}[teorema]{Corolário}

\theoremstyle{definition}
\newtheorem{definicao}[teorema]{Definição}
\newtheorem{exemplo}[teorema]{Exemplo}
\newtheorem{exercicio}{Exercício}[section]

\theoremstyle{remark}
\newtheorem{observacao}[teorema]{Observação}
\newtheorem{contraexemplo}[teorema]{Contraexemplo}
\newtheorem{intuicao}[teorema]{Intuição}

\newcommand{\R}{\mathbb{R}}
\newcommand{\N}{\mathbb{N}}
\newcommand{\Z}{\mathbb{Z}}
\newcommand{\Q}{\mathbb{Q}}
\DeclareRobustCommand{\eps}{\varepsilon}
\DeclareRobustCommand{\vet}[1]{\vec{#1}}
\DeclareRobustCommand{\dif}{\mathrm{d}}
\DeclareRobustCommand{\abs}[1]{\left\lvert #1 \right\rvert}
\DeclareRobustCommand{\norm}[1]{\left\lVert #1 \right\rVert}
\DeclareMathOperator{\grad}{grad}
\DeclareMathOperator{\Div}{div}
\DeclareMathOperator{\rot}{rot}

\newcommand{\base}[1]{\par\medskip\noindent\textbf{Base para relembrar.} #1\par\medskip}
\newcommand{\cuidado}[1]{\par\medskip\noindent\textbf{Quando a hipótese falha.} #1\par\medskip}

\newcommand{\figura}[2]{%
  \begin{figure}[htbp]
  \centering
  \IfFileExists{#1}{%
    \includegraphics[width=0.84\textwidth]{#1}%
  }{%
    \fbox{\parbox{0.8\textwidth}{\centering Figura ainda não gerada:\\ \texttt{#1}}}%
  }
  \caption{#2}
  \end{figure}%
}

\newenvironment{esqueleto}{%
  \par\medskip\noindent\textbf{Esqueleto da demonstração.}\par\nopagebreak
  \begin{enumerate}[leftmargin=1.4em]
}{\end{enumerate}\medskip}

\title{Solução --- Desigualdade triangular e Cauchy--Schwarz}
\author{}
\date{}
\begin{document}
\maketitle
\tableofcontents

As quatro demonstrações seguem a ordem do caderno de enunciados. Cada uma inclui os fatos auxiliares que ela usa.

\section{Desigualdade triangular em \texorpdfstring{$\R$}{R}}

\begin{teorema}
Para todos $x,y\in\R$, $\abs{x+y}\le\abs{x}+\abs{y}$.
\end{teorema}

\begin{proof}
Usa-se que $\R$ é corpo ordenado: a ordem é tricotômica e transitiva, $x<y$ implica $x+z<y+z$, e $x<y$ com $z>0$ implica $xz<yz$. Escreve-se $x\le y$ quando $x<y$ ou $x=y$.

Primeiro, $t^2\ge 0$ para todo $t$, com igualdade só em $t=0$. Se $t>0$, então $t\cdot t>0$. Se $t<0$, então $-t>0$ e $t^2=(-t)^2>0$. A identidade $(-a)(-a)=a^2$ vem de $a(-a)=-(a^2)$ e de $(-a)(-a)=-(a(-a))$. Além disso $0<1$: se $1<0$, então $0<-1$, e multiplicar $1<0$ por $-1>0$ daria $-1<0$, contradizendo a tricotomia.

Se $x\ge 0$ e $y\ge 0$, então $xy\ge 0$. Se $a\le b$ e $c\le d$, então $a+c\le b+d$: some $c$ à primeira desigualdade e $b$ à segunda, e use a transitividade. Se $x<y$, então $-y<-x$, somando $-x-y$. O caso não estrito segue separando igualdade e desigualdade estrita.

Agora o módulo. Se $x\ge 0$, $\abs{x}=x\ge 0$. Se $x<0$, $\abs{x}=-x>0$. Em ambos, $-\abs{x}\le x\le\abs{x}$: no caso $x\ge 0$, $-\abs{x}=-x\le 0\le x$; no caso $x<0$, $-\abs{x}=x$ e $x<0<-x$. Também $\abs{xy}=\abs{x}\abs{y}$, por inspeção dos quatro sinais, e $\abs{-x}=\abs{x}$.

Caracterização. Seja $a\ge 0$. Então $\abs{z}\le a$ se e somente se $-a\le z\le a$. Se $\abs{z}\le a$, já se tem $-\abs{z}\le z\le\abs{z}$. Se $\abs{z}=a$, então $-\abs{z}=-a$; se $\abs{z}<a$, então $-a<-\abs{z}$. Logo $-a\le z\le a$. Reciprocamente, se $-a\le z\le a$ e $z\ge 0$, então $\abs{z}=z\le a$. Se $z<0$, $\abs{z}=-z$; de $-a\le z$ segue $-z\le a$, separando igualdade e desigualdade estrita ao multiplicar por $-1$.

Aplique $-\abs{x}\le x\le\abs{x}$ e $-\abs{y}\le y\le\abs{y}$. Somando,
\[
-\bigl(\abs{x}+\abs{y}\bigr)\le x+y\le\abs{x}+\abs{y}.
\]
O número $\abs{x}+\abs{y}$ é não negativo. A caracterização com $z=x+y$ dá $\abs{x+y}\le\abs{x}+\abs{y}$.
\end{proof}

\section{Diferença dos módulos}

\begin{corolario}
Para todos $x,y\in\R$, $\abs{\,\abs{x}-\abs{y}\,}\le\abs{x-y}$.
\end{corolario}

\begin{proof}
Pela desigualdade triangular aplicada a $(x-y)+y$,
\[
\abs{x}\le\abs{x-y}+\abs{y},
\]
logo $\abs{x}-\abs{y}\le\abs{x-y}$. Trocando $x$ e $y$, $\abs{y}-\abs{x}\le\abs{y-x}$. Como $\abs{y-x}=\abs{-(x-y)}=\abs{x-y}$, tem-se $\abs{y}-\abs{x}\le\abs{x-y}$. Multiplicar por $-1$ inverte o sentido, com o cuidado de que a igualdade permanece igualdade. Assim
\[
-\abs{x-y}\le\abs{x}-\abs{y}\le\abs{x-y}.
\]
A caracterização do módulo, com o número não negativo $\abs{x-y}$, dá $\abs{\,\abs{x}-\abs{y}\,}\le\abs{x-y}$.
\end{proof}

\section{Cauchy--Schwarz em \texorpdfstring{$\R^n$}{Rn}}

\begin{teorema}
Para todos $\vet{u},\vet{v}\in\R^n$, $\abs{\vet{u}\cdot\vet{v}}\le\norm{\vet{u}}\norm{\vet{v}}$.
\end{teorema}

\begin{proof}
O produto escalar é bilinear e simétrico, porque a soma e o produto de $\R$ o são em cada coordenada. Em particular
\[
\norm{\vet{u}+t\vet{v}}^2
=\norm{\vet{u}}^2+2t\,(\vet{u}\cdot\vet{v})+t^2\norm{\vet{v}}^2
\]
para todo $t\in\R$, e o primeiro membro é não negativo. Também $\norm{\vet{v}}^2=\sum v_i^2=0$ se e somente se $\vet{v}=\vet{0}$. Se $b\ge 0$, $c\ge 0$ e $b^2=c^2$, então $b=c$, porque $(b-c)(b+c)=0$ e $b=-c$ com ambos não negativos força $b=c=0$. Daí $\sqrt{w^2}=\abs{w}$. Se $0\le\alpha\le\beta$, então $\sqrt{\alpha}\le\sqrt{\beta}$. Caso contrário $\sqrt{\beta}<\sqrt{\alpha}$. As duas raízes são não negativas, e
\[
\alpha-\beta
=\bigl(\sqrt{\alpha}-\sqrt{\beta}\bigr)\bigl(\sqrt{\alpha}+\sqrt{\beta}\bigr)>0,
\]
porque os dois fatores são positivos. Isso contradiz $\alpha\le\beta$.

Se $\vet{v}=\vet{0}$, os dois lados da desigualdade pedida são zero.

Se $\vet{v}\neq\vet{0}$, seja $q(t)=At^2+Bt+C$ com $A=\norm{\vet{v}}^2>0$, $B=2(\vet{u}\cdot\vet{v})$ e $C=\norm{\vet{u}}^2$. Então $q(t)\ge 0$ para todo real $t$. O discriminante $\Delta=B^2-4AC$ não pode ser positivo. Se fosse, as raízes $t_{\pm}=(-B\pm\sqrt{\Delta})/(2A)$ satisfariam $t_-<t_+$ e
\[
t_-+t_+=-\frac{B}{A},\qquad
t_-t_+=\frac{B^2-\Delta}{4A^2}=\frac{C}{A},
\]
porque $B^2-\Delta=4AC$. Logo $q(t)=A(t-t_-)(t-t_+)$. Entre as raízes o produto $(t-t_-)(t-t_+)$ é negativo e $A>0$, então $q(t)<0$, absurdo. Portanto $\Delta\le 0$, isto é
\[
(\vet{u}\cdot\vet{v})^2\le\norm{\vet{u}}^2\norm{\vet{v}}^2.
\]
Os dois lados são não negativos. A raiz quadrada preserva a ordem, e $\sqrt{(\vet{u}\cdot\vet{v})^2}=\abs{\vet{u}\cdot\vet{v}}$, enquanto $\sqrt{\norm{\vet{u}}^2\norm{\vet{v}}^2}=\norm{\vet{u}}\norm{\vet{v}}$.
\end{proof}

\section{Desigualdade triangular da norma euclidiana}

\begin{teorema}
Para todos $\vet{u},\vet{v}\in\R^n$, $\norm{\vet{u}+\vet{v}}\le\norm{\vet{u}}+\norm{\vet{v}}$.
\end{teorema}

\begin{proof}
Expanda o quadrado da norma pela bilinearidade:
\[
\norm{\vet{u}+\vet{v}}^2
=\norm{\vet{u}}^2+2\,\vet{u}\cdot\vet{v}+\norm{\vet{v}}^2.
\]
Para todo real $t$, $t\le\abs{t}$: se $t\ge 0$ há igualdade, e se $t<0$ então $t<0<\abs{t}$. Logo $\vet{u}\cdot\vet{v}\le\abs{\vet{u}\cdot\vet{v}}$ e
\[
\norm{\vet{u}+\vet{v}}^2
\le\norm{\vet{u}}^2+2\abs{\vet{u}\cdot\vet{v}}+\norm{\vet{v}}^2.
\]
Cauchy--Schwarz substitui o módulo do produto escalar:
\[
\norm{\vet{u}+\vet{v}}^2
\le\norm{\vet{u}}^2+2\norm{\vet{u}}\norm{\vet{v}}+\norm{\vet{v}}^2
=\bigl(\norm{\vet{u}}+\norm{\vet{v}}\bigr)^2.
\]
Os números $\norm{\vet{u}+\vet{v}}$ e $\norm{\vet{u}}+\norm{\vet{v}}$ são não negativos, e a desigualdade dos quadrados passa para as raízes, porque a raiz quadrada é crescente em $[0,\infty)$, como registrado na prova anterior. Portanto $\norm{\vet{u}+\vet{v}}\le\norm{\vet{u}}+\norm{\vet{v}}$.
\end{proof}

\end{document}
